\section{Specializing the logarithmic square}
\label{sec:square-specialization}

We now isolate the numerical consequence of the extra-component
vanishing.  All intersections in this section are ordinary projective
intersection numbers; the conversion to the local invariant is exactly
Lemma~\ref{lem:projective-local-square}.

\begin{lemma}[Square specialization]
\label{lem:square-specialization}
Let $q:Y\to C$ be a flat projective morphism to a smooth connected
complex curve, with pure two-dimensional fibres.  Let $L$ be a
rational Cartier divisor on $Y$.  If
$[Y_0]=\sum_j m_j[S_j]$ is the cycle of a fibre, then for every $c\in C$
\[
 (L|_{Y_c})^2=\sum_j m_j(L|_{S_j})^2.
\]
In particular, the multiplicities on the right are one when $Y_0$ is
reduced.
\end{lemma}

\begin{proof}
Choose $r>0$ such that $rL$ is Cartier.  Flatness and properness imply
that $\chi(Y_c,\mathcal O_Y(nrL)|_{Y_c})$ is constant in $c$ for every
integer $n$ \cite[Lemma~36.32.2]{StacksProject}.  This is a polynomial
in $n$ of degree at most two, whose quadratic coefficient is
$r^2(L|_{Y_c})^2/2$.  Its leading term is the sum of the leading terms
on the two-dimensional components, weighted by their generic lengths
\cite[Lemmas~33.45.1, 33.45.2 and 33.45.6]{StacksProject}.  Comparing
these coefficients and dividing by $r^2$ proves the assertion.
\end{proof}

\begin{proposition}[The numerical conclusion of the log model]
\label{prop:beta-from-log-model}
Let $X\to C$, with section $\Gamma$, be the projective family of
Lemma~\ref{lem:compactification}.  Let $h:Y\to X$ be a projective
birational morphism from a normal integral threefold, isomorphic off
$\Gamma$.  Suppose $Y_0=S_0+\sum_a S_a$ is reduced Cartier, $S_0$
maps birationally to $X_0$, and the $S_a$ map to $\Gamma(0)$.  Suppose
that, after shrinking $C$, the other fibres $Y_c$ are normal and map
birationally to $X_c$.

Let $H$ be the reduced horizontal exceptional boundary and set
$L=K_Y+H$, assumed rational Cartier and $h$-nef.  For each component
$S$ of $Y_0$, let $B_S$ be the reduced divisor of its intersections
with the other components of $Y_0+H$.  Assume that $(S,B_S)$ is log
canonical and
\[
 L|_S\sim_{\mathbb Q} K_S+B_S.
\]
For $S_0$, require that $B_{S_0}$ is the full reduced inverse image of
$\Gamma(0)$.  For $c\ne0$, require likewise that the full reduced
inverse image $B_c$ of $\Gamma(c)$ is a divisor, that $(Y_c,B_c)$ is
log canonical, and that
$L|_{Y_c}\sim_{\mathbb Q}K_{Y_c}+B_c$.
If
\[
 (K_{S_a}+B_{S_a})^2=0\qquad\text{for every }a,
\]
then $\beta(X_c,\Gamma(c))=\beta(X_0,\Gamma(0))$ for all remaining
$c\in C$.
\end{proposition}

\begin{proof}
The integral threefold $Y$ dominates the smooth curve, so it is flat:
its local rings are torsion-free modules over the corresponding
discrete valuation rings.  Its fibres have pure dimension two, and
it is projective over $C$.  Lemma~\ref{lem:square-specialization} and
the assumed vanishing give, for $c\ne0$,
\[
 (K_{Y_c}+B_c)^2
 = (L|_{Y_c})^2
 = (L|_{S_0})^2+\sum_a(L|_{S_a})^2
 = (K_{S_0}+B_{S_0})^2.
\]
The restrictions of $L$ are nef over the respective $X_c$, since
$L$ is $h$-nef.  The full-boundary hypotheses therefore allow
Lemma~\ref{lem:projective-local-square} on both the central and
noncentral birational models.  The displayed equality and
$K_{X_c}^2=K_{X_0}^2$ yield the claimed equality of local invariants.
\end{proof}

When a finite change of curve was used to construct $Y$, this equality
still concerns the original fibre singularities: such a change
repeats individual fibres and does not alter their analytic germs.

Lemma~\ref{lem:supported-square} extends the specialization identity to
nonproper families when the divisor has proper support over the curve.
Its two proofs explain how the numerical comparison can be kept local.
